% MT035 / explanation — canonical editable content.
先把概率方向写清：设\(p=P(\text{垃圾邮件}\mid x)\)，当前规则是\(p>t\)就拦截。用户“几乎收不到邮件”意味着大量正常邮件被误判为垃圾，这类错误代价很高。


\textbf{训练层面的修复：提高正常邮件的权重。} 数据中垃圾邮件很多，普通平均损失可能过度迎合多数类。给正常邮件更高的权重，使误拦一封正常邮件付出更多代价，促使模型学习区分它们的特征。补充更多有代表性的正常邮件也是可选的数据改进。


\textbf{决策层面的修复：提高判垃圾的阈值。} 例如原来\(t=0.5\)，某封正常邮件被估计垃圾概率为0.6就被拦截；把教学示例阈值升到0.8后，它会放行。这样需要更强的垃圾证据才拦截，代价是可能漏过更多真正垃圾邮件。


若模型输出的是“正常邮件概率”，同一个偏好应表现为降低判正常的门槛。所以不要脱离标签背“阈值要提高”或“要降低”，先写出哪个分数大时判哪一类。


在验证集上选择权重与阈值，并记录正常邮件被误拦的比例及垃圾漏过情况。


\textbf{出处：} 2021B*期中 Q9，原卷第4页（10分）


\textbf{答案依据：} 2021B*期中 Q9 随卷／配套参考答案；条件与解释经复核。
